% ===== §11 =====
\section{Ethical Legitimacy Observed across Traditions}
\label{sec:traditions}

\noindent
The legitimacy that separates understanding from its counterfeit has been stated through non-possession, and the previous section closed by holding that a single ethical rendering should not carry the whole of it. This section draws the ground more fully. It observes the legitimacy of understanding across several ethical traditions, each of which renders an aspect of what it is to treat another rightly and leaves an aspect for the others to render, and it treats non-possession as one such aspect among them. The traditions are taken in their secular and argued forms, as accounts of right conduct that stand on reasons open to examination, and the section passes through them in the manner the epistemology prescribes, to locate in each what it reaches and what it leaves.

The functional case already brought one tradition into play, though it did so to mark its insufficiency. A consequentialist ethic renders the rightness of an act in its results, the good it produces and the harm it forestalls, and on such an ethic the legitimacy of an understanding would lie in the goods it brings to the bond, the trust, the shared value, the eased conflict, the sustained quality of a life. Part Three showed why this rendering, taken alone, cannot draw the distinction the account requires, since a refined possession of the other can produce these goods, for a time, as efficiently as understanding. The consequentialist rendering reaches the real and important truth that a legitimate understanding is fruitful, that it does good in the bond, and it leaves untouched the matter of how the other is treated in the producing of that good, which is the matter on which understanding and its counterfeit divide. Its aspect is the fruitfulness of understanding; its omission is the standing of the other.

A deontological ethic in the tradition of Kant~\citep{kant1785groundwork} renders right conduct through duty and the moral law, and it supplies, at its center, a principle that speaks directly to the present matter, that one is to treat humanity, in oneself and in others, never merely as a means and always at the same time as an end. This principle gives the standing of the other the consequentialist rendering omits, and it gives it as a requirement that no calculation of goods can override. The counterfeit of understanding, which takes the other into a system of management, treats the other merely as a means, and the Kantian rendering names this the precise wrong that the consequentialist accounting cannot register. Its aspect is the dignity of the other as an end, a standing that the possessive relation violates however favorable its results. What it leaves is the concreteness and the relationality of the bond, for the Kantian ground is a universal law addressed to any rational being as such, and it renders the other as a bearer of a dignity shared by all rather than as this particular other whose singular world is to be entered. It secures that the other is not used, and it says little of what it is to attend to just this person in their difference.

A care ethic, in the work of Noddings and of Gilligan~\citep{noddings1984caring,gilligan1982voice}, renders right conduct through the attentive response to a particular other in their need, and it supplies exactly the concreteness the Kantian ground leaves abstract. On this rendering legitimacy is a matter of attentiveness, of receiving the other in their particularity and responding to what they, and not humanity in general, require. It gives the positive and relational content that non-possession, stated as a restraint, does not carry on its own, for non-possession says what a legitimate understanding refrains from, the annexation of the other, and the care ethic says what it actively does, the attentive entry into and response to the other's world. Its aspect is the responsive attention to the singular other. What it leaves is the matter of limit and of mutuality, for an ethic of care, pressed to its edge, can pass into a self-effacement in which the one who cares is drawn without return, and this is a relational imbalance of its own, a failure of the mutual generation the account requires, which the care rendering is not well placed to mark.

An ethic of virtue in the tradition of Aristotle~\citep{aristotle2000ethics} renders right conduct through the character of the agent and the flourishing it makes possible, and in its account of friendship it reaches the matter of mutuality the care ethic leaves. The friendship of the best kind, on this account, is one in which each wishes the good of the other for the other's own sake, and in which the two flourish together in a shared activity, each the occasion of the other's excellence. This renders legitimacy as a mutual and generative matter, a good realized between two who are equally its authors, and it supplies the reciprocity that the asymmetry of care and the universality of duty do not frame. Its aspect is the mutual flourishing of two who wish each other's good for its own sake, which is close to the account's own conception of a bond in which two worlds generate what neither held alone. What it leaves is its reliance on a substantive account of flourishing, a settled conception of the human good toward which the friendship tends, which sits in tension with the account's refusal of a predetermined end, and this tension is itself a matter the observation must hold rather than resolve.

Set these renderings into relation, and the legitimacy of understanding becomes observable in a way no one of them affords. The consequentialist reaches its fruitfulness and omits the standing of the other; the Kantian reaches that standing as dignity and omits the concrete singular; the care ethic reaches the responsive attention to the singular and omits the limit and the mutuality; the virtue ethic reaches the mutual flourishing and relies on a settled good the account cannot grant. Non-possession stands among these as one aspect, nearest to the Kantian refusal of the other as mere means, and it is supplied, by the others, with the positive attention, the reciprocity, and the fruitfulness that a restraint stated alone does not carry. What legitimacy is, in the treatment of the other that separates understanding from extraction, is observed in the relations among these renderings, as the account holds every concept to be observed, and it is exhausted by none of them.

\begin{claim}[Legitimacy observed across traditions]
The ethical legitimacy that separates understanding from its counterfeit is a concept, and by the paper's own method it is exhausted by no single ethical tradition and becomes observable through the relations among several. A consequentialist rendering reaches the fruitfulness of understanding; a Kantian rendering reaches the standing of the other as an end and not a mere means; a care ethic reaches the attentive response to the singular other; an Aristotelian ethic of friendship reaches the mutual flourishing of two who wish each other's good for its own sake. Each omits what the others supply. Non-possession is one aspect among these, nearest the Kantian refusal of the other as means, and completed by the positive attention, the reciprocity, and the fruitfulness the others render. Legitimacy is observed in the relation among the traditions and is exhausted by none.
\end{claim}

This section has done, at the level of ethics, what the paper argues understanding always does. It has taken a concept, the legitimacy of understanding, that no single symbolic system renders without remainder, and it has brought it into view through the relations among several renderings of it, each partial, none final. The normative ground of the paper is thereby drawn in the same manner as its epistemology, and the earlier reliance on non-possession is corrected by observing non-possession as one aspect within a relation of traditions that together render more of the legitimacy than any renders alone, and not by replacing one sole criterion with another. The normative case is, in its own construction, an instance of the relational understanding it concerns.

Two conditions hold this observation to the paper's stance, as they held the account of legitimacy in the previous section. The first is that observing legitimacy across traditions does not yield a correct ethical criterion, a formula that would settle, for every case, whether an understanding is legitimate. It yields an observation, partial and open, of what legitimacy involves, in which the traditions check and supplement one another without converging on a single rule. The account claims no more ethical correctness than it claims epistemic correctness, and for the same reason. The second is that the field of traditions is itself historical. These renderings arose in the development of moral life and will be joined by others as that life develops, and the observation of legitimacy they afford is generated in history and open to its further generation, set as no final criterion. The account offers this relation of traditions as the observation of legitimacy that the present development of ethical life affords, held as the best the present yields and answerable to what comes after, in the manner the framework holds of all its commitments.

With this the normative case is complete. Function tells what understanding does; legitimacy, observed across the traditions that render it, tells understanding from its counterfeit, without a correct criterion and within the history that generates it. The paper turns now to the problem that sets the cultivation of these capacities in intimacy apart from their cultivation anywhere else.
